% ============================================================
% example-paper-twocolumn.tex — harryopo-paper.cls 双栏模式完整示例
% 展示：跨栏标题摘要、跨栏图表、算法环境、栏平衡、上标引用
% ============================================================
\documentclass[twocolumn]{harryopo-paper}

% --- 论文标题 ---
\title{边缘计算环境下基于Flink的实时数据流处理与动态资源调度优化}

% --- 作者信息 ---
\author{王五\textsuperscript{1} \quad 赵六\textsuperscript{2}\\%
       {\small\textsuperscript{1} 计算机学院，深圳 518172}\\
       {\small\textsuperscript{2} 哈尔滨工业大学(深圳) 计算机科学与技术学院，深圳 518055}}
\date{2026年6月21日}

% --- 摘要存储（用于 \maketitlewithabstract） ---
\abstractcontent{随着物联网设备的快速普及和海量实时数据的产生，传统云计算架构在处理边缘侧实时数据时面临高延迟、带宽瓶颈和隐私保护不足等挑战。本文针对边缘计算环境下的实时数据流处理问题，提出了一种基于Apache Flink的分布式流处理框架，并设计了动态资源调度优化策略。实验结果表明，所提框架在毫秒级延迟约束下，吞吐量较基准方法提升了32\%，资源利用率提高了28\%。}

\keywordscontent{边缘计算；实时流处理；Apache Flink；动态资源调度}

\begin{document}

% 双栏模式下使用 maketitlewithabstract 实现标题+摘要跨栏排版
\maketitlewithabstract

\section{引言}

边缘计算作为一种新型计算范式，通过将计算能力从云端下沉到网络边缘，为终端设备提供就近服务\cite{edge}。随着5G技术的商用化推进，边缘计算在智能制造、自动驾驶、智慧医疗等领域的应用前景愈发广阔。

实时数据流处理是边缘计算的核心需求之一。与传统批处理不同，流处理要求系统在毫秒级延迟内完成数据采集、过滤、聚合和响应。Apache Flink\cite{flink}作为一个真正的流处理引擎，以其精确一次（Exactly-once）语义和低延迟特性，成为边缘计算场景下的首选流处理框架。

\section{系统架构}

\subsection{云-边-端三层架构}

本文采用"云-边-端"三层协作架构。云端负责模型训练、全局调度和历史数据存储；边缘层部署Flink集群，负责实时流处理和本地决策；终端层负责数据采集和指令执行。

\subsection{Flink流处理引擎}

在边缘节点上，我们部署了轻量化的Flink实例，配置了基于RocksDB的状态后端以支持大状态管理，并启用了检查点（Checkpoint）机制保障容错性。

\subhead{状态后端配置}RocksDB状态后端能将状态数据持久化到本地磁盘，适合大状态场景。在边缘节点内存有限的条件下，这一特性尤为重要。

\subhead{检查点优化}采用增量检查点策略，仅保存自上次检查点以来的状态变更，大幅减少了检查点开销。

\section{动态资源调度策略}

\subsection{问题建模}

边缘节点的计算资源有限且动态变化，因此资源调度问题可以建模为约束优化问题。设边缘节点集合为$\mathcal{N} = \{n_1, n_2, \dots, n_k\}$，任务集合为$\mathcal{T} = \{t_1, t_2, \dots, t_m\}$。

目标函数为最小化加权延迟与资源成本的组合：
\begin{equation}
\min \sum_{i=1}^{m} \sum_{j=1}^{k} \bigl(w_1 \cdot x_{ij} \cdot L(t_i, n_j) + w_2 \cdot x_{ij} \cdot C(t_i, n_j)\bigr)
\label{eq:objective}
\end{equation}

约束条件包括：
\begin{itemize}
    \item 每个任务必须分配到恰好一个节点：$\sum_{j} x_{ij} = 1, \forall i$
    \item 节点资源上限约束：$\sum_{i} x_{ij} \cdot r(t_i) \leq R(n_j), \forall j$
    \item 延迟SLA约束：$L(t_i, n_j) \leq L_{\max}, \forall (i,j): x_{ij}=1$
\end{itemize}

\subsection{算法设计}

采用改进的启发式贪心算法进行在线调度，如算法\ref{alg:scheduling}所示。

\begin{algorithm}[t]
\caption{在线动态资源调度算法}
\label{alg:scheduling}
\KwIn{任务队列 $\mathcal{Q}$，节点列表 $\mathcal{N}$，SLA约束 $L_{\max}$}
\KwOut{任务-节点映射 $\mathcal{M}$}
\SetKwFunction{Score}{Score}

\BlankLine
\ForEach{$t \in \mathcal{Q}$}{
    $\mathcal{C} \leftarrow \{n \in \mathcal{N} \mid \text{满足资源约束}\}$\;
    \If{$\mathcal{C} = \emptyset$}{
        将 $t$ 加入等待队列\;
        \textbf{continue}\;
    }
    $n^* \leftarrow \arg\min_{n \in \mathcal{C}} \Score(t, n)$\;
    \If{$\Score(t, n^*) \leq L_{\max}$}{
        $\mathcal{M}[t] \leftarrow n^*$\;
        更新 $n^*$ 的资源状态\;
    }
    \Else{
        将 $t$ 升级到云端处理\;
    }
}
\Return{$\mathcal{M}$}
\end{algorithm}

\section{实验评估}

\subsection{实验环境}

实验使用3台边缘服务器（Intel Xeon E-2278G, 32GB RAM）和10台模拟终端设备。Flink集群配置为3个TaskManager，每个分配8GB堆内存。

\subsection{吞吐量对比}

表\ref{tab:throughput}展示了不同框架在相同工作负载下的吞吐量对比。

\begin{table}[t]
\centering
\small
\begin{tabular}{lcc}
\toprule
\textbf{框架} & \textbf{低负载} & \textbf{高负载} \\
\midrule
Apache Storm & 12,500 & 38,200 \\
Spark Streaming & 9,800 & 45,600 \\
\textbf{本文方法} & \textbf{18,200} & \textbf{60,800} \\
\bottomrule
\end{tabular}
\caption{流处理框架吞吐量对比（条/秒）}
\label{tab:throughput}
\end{table}

\subsection{延迟分析}

在低负载（1000条/秒）条件下，本文方法的P99延迟为8.2ms，显著优于Storm的15.4ms和Spark Streaming的120ms（微批次导致）。在高负载（5000条/秒）条件下，本文方法通过动态调度策略仍将P99延迟控制在23.5ms以内。

\subsection{资源利用率}

动态调度策略使CPU利用率从静态分配方案的62\%提升至78\%，内存利用率从55\%提升至72\%。

% ==== 跨栏大图 ====
\begin{figure*}[t]
\centering
\colorbox{gray!8}{\parbox{0.9\textwidth}{\centering\vspace{3cm}
    {\Large\color{gray} [跨双栏系统架构全图]}\\[0.5cm]
    \small 云端（模型训练/全局调度） $\leftrightarrow$ 边缘层（Flink集群/实时推理） $\leftrightarrow$ 终端层（传感器/执行器）
\vspace{3cm}}}
\caption{云-边-端三层系统总体架构（跨栏排版）}
\label{fig:architecture_full}
\end{figure*}

\section{相关工作}

边缘计算的概念最早由Shi等人\cite{edge}系统阐述。在流处理方面，Carbone等人\cite{flink}提出的Apache Flink以其真正的流处理语义成为业界标准。近年来，研究者们提出了多种边缘-云端协同的流处理方案。

与现有工作相比，本文的主要区别在于：(1) 针对边缘资源受限场景设计了轻量化Flink配置；(2) 提出了在线动态资源调度算法，而非静态分配；(3) 在真实边缘服务器环境下进行了充分验证。

\section{结论与展望}

本文提出了一种基于Apache Flink的边缘实时数据流处理框架，并设计了动态资源调度优化策略。实验结果表明，所提方法在吞吐量和资源利用率方面均有显著提升。

未来的研究方向包括：探索基于强化学习的智能调度策略、支持异构硬件（GPU/FPGA）的流处理加速、以及在更大规模边缘集群上进行验证。

% ==== 参考文献 ====
\begin{thebibliography}{9}
\bibitem{edge} Shi W, Cao J, Zhang Q, et al. Edge computing: Vision and challenges[J]. IEEE Internet of Things Journal, 2016, 3(5): 637-646.
\bibitem{flink} Carbone P, Katsifodimos A, Ewen S, et al. Apache Flink: Stream and batch processing in a single engine[J]. Bulletin of the IEEE Computer Society Technical Committee on Data Engineering, 2015, 38(4): 28-38.
\bibitem{storm} Toshniwal A, Taneja S, Shukla A, et al. Storm@Twitter[C]. SIGMOD, 2014: 147-156.
\bibitem{spark} Zaharia M, Xin R S, Wendell P, et al. Apache Spark: A unified engine for big data processing[J]. Communications of the ACM, 2016, 59(11): 56-65.
\end{thebibliography}

\end{document}
